% PDF-to-TeX transcription by the assistant; not the original downloaded file.
\begin{theorem}[1.1 (DiGeorgi--Nash--Moser)]
Let
\[Lu:=\sum\partial_i(a_{ij}\partial_j u)\quad\text{and}\quad0<\lambda\leq a_{ij}\leq\Lambda.\tag{DGH}\]
Then there exists $\alpha(n,\lambda,\Lambda)>0$ and $C(n,\lambda,\Lambda)$ such that if $Lu=0$, then
\[\|u\|_{C^\alpha(B_{1/2})}\leq C(\lambda,\Lambda,n)\|u\|_{C^0(B_1)}.\]
\end{theorem}
Note that this estimate does not in any way involve derivatives of the $a_{ij}$.
\begin{proposition}[1.3]
If $L$ follows (DGH) and $Lu=0$ on $B_1$ then
\[\int_{B_{1/2}}|\nabla u|^2\lesssim\int_{B_1}|u|^2.\]
\end{proposition}
\begin{proof}
We will use integration by parts and localization. Let $\eta=1$ on $B_{1/2}$ and be $0$ outside of $B_1$.
\begin{align*}
\int_{B_{1/2}}|\nabla u|^2&\leq\int\eta^2|\nabla u|^2\approx\int\eta^2\sum a_{ij}\partial_i u\partial_j u\\
&\leq\int\eta^2(Lu)u+\int|\nabla\eta|\eta|\nabla u||u|\\
&\leq\left(\int\eta^2|\nabla u|^2\right)^{1/2}\left(\int|\nabla\eta|^2u^2\right)^{1/2}.
\end{align*}
\end{proof}
\paragraph{2. $L^\infty$ Bound}
\begin{theorem}[2.1 (DGNM $L^\infty$ bound)]
Let $L$ satisfy (DGH), $Lu\geq0$, $u>0$. Then
\[\|u\|_{L^\infty(B_{1/2})}\leq\|u\|_{L^2(B_1)}.\]
\end{theorem}
We start with a lemma:
\begin{lemma}[2.1]
Under the hypotheses, and if $1/2\leq r<r+w\leq1$ then
\[\|\nabla u\|_{L^2(B_r)}\lesssim\|u\|_{L^2(B_{r+w})}w^{-1}.\]
\end{lemma}
\begin{proof}
Let $\eta=1$ on $B_r$ and $0$ on $B_{r+w}^c$. Note that $\eta$ can be constructed so that $|\nabla\eta|<2w^{-1}$. Then the proof proceeds in exactly the same fashion as Proposition 1.3.
\end{proof}
\begin{lemma}[2.2]
Under hypotheses, and $1/2\leq r<r+2\leq1$, we have
\[\|u\|_{L^{2n/(n-2)}(B_r)}\lesssim w^{-1}\|u\|_{L^2(B_{r+w})}.\]
\end{lemma}
\begin{proof}
Consider $\eta u$ with $\eta=1$ on $B_r$, and $0$ outside of $B_{r+w/2}$. Then by the Sobolev inequality, we have
\[\|\eta u\|_{L^{2n/(n-2)}}\lesssim\|\nabla(\eta u)\|_{L^2}\leq\|(\nabla\eta)u\|_{L^2}+\|\eta(\nabla u)\|_{L^2}.\]
Also, we have that
\begin{align*}
\|(\nabla\eta)u\|_{L^2}&\leq\|\nabla\eta\|_\infty\|u\|_{L^2(B_{r+w/2})}\lesssim w^{-1}\|u\|_{L^2(B_{r+w})},\\
\|\eta(\nabla u)\|_{L^2}&\leq\|\nabla u\|_{L^2(B_{r+w/2})}\lesssim w^{-1}\|u\|_{L^2(B_{r+w})}.
\end{align*}
\end{proof}
\begin{lemma}[2.3]
If $\beta>1$, $Lu\geq0$ and $u>0$, then $Lu^\beta\geq0$.
\end{lemma}
\begin{proof}
Compute:
\begin{align*}
Lu^\beta&=\sum\partial_i(a_{ij}\partial_j(u^\beta))=\sum\partial_i(a_{ij}\beta u^{\beta-1}\partial_j u)\\
&=(Lu)(\beta u^{\beta-1})+\sum a_{ij}\partial_i u\partial_j u\,\beta(\beta-1)u^{\beta-2}\geq0
\end{align*}
where the last inequality comes from ellipticity of $a_{ij}$.
\end{proof}
Now, apply Lemma 2.2 to $u^\beta$ to get
\[\|u^\beta\|_{L^{2n/(n-2)}(B_r)}\lesssim w^{-1}\|u^\beta\|_{L^2(B_{r+w})}.\]
Rewriting this with $s=n/(n-2)$ we get
\begin{lemma}[2.4]
If $1/2\leq r<r+w\leq1$ and $p\geq2$, then
\[\|u\|_{L^{sp}(B_r)}\leq(Cw^{-1})^{2/p}\|u\|_{L^p(B_{r+w})}.\]
\end{lemma}
For the next step, we iterate this lemma. If we have $1=r_0>r_1>\cdots>r_k>1/2$, then we get the sequence of inequalities
\[\|u\|_{L^2(B_1)}\geq A_0\|u\|_{L^{2s}(B_{r_1})}\geq\cdots\geq A_0\cdots A_{k-1}\|u\|_{L^{2s^k}(B_{r_k})}\]
where the $A_j$ are given by Lemma 2.4. Let's pick $r_j=1/2+1/(j+2)$, so that $r_j-r_{j+1}\approx j^{-2}$. Thus, $A_j=(C(r_j-r_{j-1})^{-1})^{s^{-j}}$. Therefore,
\begin{align*}
\log\Big(\prod A_j\Big)&\leq\sum\log(A_j)\leq\sum_{j=0}^\infty s^{-j}\big(C+C\log(r_j-r_{j+1})\big)\\
&\leq\sum_{j=0}^\infty s^{-j}(C+C\log j)<\infty.
\end{align*}
\paragraph{3. Finishing the Proof}
We will show a Harnack inequality for our $L$ which satisfies (DGH).
\begin{theorem}[3.2 (DGNM Harnack)]
If $L$ satisfies (DGH), $Lu=0$, $1>u>0$ on $B_1$, and
\[|\{x\in B_{1/2}:u(x)>1/10\}|\geq\frac1{10}|B_{1/2}|,\tag{P}\]
then $\min_{B_{1/2}}u\geq\gamma(n)$.
\end{theorem}
For now, let's assume this theorem, and see how it implies the DiGeorgi-Nash-Moser estimate.
\begin{definition}[3.1]
$\operatorname{osc}_\Omega u:=\sup_\Omega u-\inf_\Omega u$.
\end{definition}
\begin{corollary}[3.1]
If $Lu=0$ on $\Omega$, $B_r(x)\subset\Omega$, then
\[\operatorname{osc}_{B_{r/2}(x)}u\leq(1-\gamma)\operatorname{osc}_{B_r(x)}u.\tag{O}\]
\end{corollary}
\begin{proof}
We start with some simple reductions via scaling. Without loss of generality, we can take:
\[\inf_{B_r(x)}u=0,\quad\sup_{B_r(x)}u=1,\quad r=1,\]
\[|\{x\in B_{1/2}:u(x)\geq1/2\}|\geq B_{1/2}/2.\]
Thus by DGNM Harnack, $\min_{B_{1/2}}u\geq\gamma$, and thus $\operatorname{osc}_{B_{1/2}}u\leq1-\gamma=(1-\gamma)\operatorname{osc}_{B_1}u$.
\end{proof}
Now we can complete the proof with the following:
\begin{proposition}[3.1]
Let $u:B_1\to\R$ satisfy (O). Then $\|u\|_{C^\alpha(B_{1/2})}\lesssim\|u\|_{C^0(B_1)}$ for some $\alpha=\alpha(\gamma)>0$.
\end{proposition}
\begin{proof}
Let $x,y\in B^{1/2}$, $|x-y|=d$ and $a=(x+y)/2$. Then
\[|u(x)-u(y)|\leq(\operatorname{osc}_{B_d(a)}u)(1-\gamma)\leq\cdots\leq(1-\gamma)^k\operatorname{osc}_{B_{2^kd}(a)}u.\]
Choose $k$ such that $1/4<2^kd\leq1/2$. Then $k=\log_2(1/d)+O(1)$, and so
\[|u(x)-u(y)|\leq(1-\gamma)^k\operatorname{osc}_{B_1}u\leq2(1-\gamma)^k\|u\|_{C^0(B_1)}.\]
Also,
\[(1-\gamma)^k\leq4(1-\gamma)^{\log_2(1/d)}=4d^{-\log_2(1-\gamma)}.\]
Therefore, setting $\alpha(\gamma)=-\log_2(1-\gamma)\approx\gamma+O(\gamma^2)$, we get our proposition.
\end{proof}
\paragraph{DGNM Harnack: the deferred proof}
\begin{lemma}[3.2]
If $L$ satisfies (DGH), $Lu=0$, $u>0$ on $B_1$ then $\|\nabla\log u\|_{L^2(B_{1/2})}\lesssim1$.
\end{lemma}
\begin{proof}
Pick a nice cutoff function $\eta$ as usual.
\begin{align*}
\int_{B_{1/2}}|\nabla\log u|^2&=\int\eta^2|\nabla\log u|^2\lesssim\int\eta^2\sum a_{ij}\partial_i\log u\partial_j\log u\\
&=\int\eta^2\sum a_{ij}\frac{\partial_i u}{u}\frac{\partial_j u}{u}
=-\int\eta^2\sum a_{ij}\partial_i u\partial_j u^{-1}\\
&\lesssim\int\eta|\nabla\eta||\nabla u|u^{-1}=\int\eta|\nabla\eta||\nabla\log u|\\
&\leq\left(\int\eta^2|\nabla\log u|^2\right)^{1/2}\left(\int|\nabla\eta|^2\right)^{1/2}.
\end{align*}
\end{proof}
Letting $w=-\log u$, we have that $\|\nabla w\|_{L^2(B_{9/10})}\lesssim1$. We want an $L^\infty$ bound on $w$. By (P), we have that
\[|\{x\in B_{1/2}:w\leq\log10\}|\geq\frac1{10}|B_{1/2}|.\]
Now we use the Poincare Inequality:
\begin{theorem}[3.3 (Poincare)]
If (P) then $\int_{B_{8/10}}|w|^2\lesssim\int_{B_{9/10}}|\nabla w|^2+1$.
\end{theorem}
Therefore, we have an $L^2$ bound on $w$ instead of $\nabla w$. Now we have
\begin{lemma}[3.3]
$Lw\geq0$.
\end{lemma}
\begin{proof}
Compute:
\begin{align*}
-\sum\partial_i(a_{ij}\partial_j\log u)&=-\sum\partial_i(a_{ij}(\partial_j u)u^{-1})\\
&=Lu\cdot u^{-1}+\sum a_{ij}(\partial_i u)(\partial_j u)u^{-2}\geq0.
\end{align*}
\end{proof}
Finally, $w=-\log u>0$ because $u<1$, and so we can apply Theorem 2.1 and get
\[\|w\|_{L^\infty(B_{1/2})}\lesssim\|w\|_{L^2(B_{8/10})}\lesssim1\]
thus completing the proof of the Harnack inequality.
